% Part: model-theory
% Chapter: models-of-arithmetic
% Section: computable-models

\documentclass[../../../include/open-logic-section]{subfiles}

\begin{document}

\olfileid[fr]{mod}{mar}{cmp}
\section{Modèles calculables de l'arithmétique}

\begin{explain}
Le modèle standard~$\Struct{N}$ possède
deux propriétés intéressantes. Son domaine
est l'ensemble~$\Nat$ des entiers naturels :
ses éléments sont précisément les objets
dont traite le langage de l'arithmétique,
et le numéral standard~$\num{n}$ y désigne
effectivement~$n$. De plus, les
interprétations de tous les symboles
non logiques de~$\Lang{L_A}$ sont
\emph{calculables}. Les fonctions
successeur, addition et multiplication
qui interprètent $\prime$, $+$ et
$\times$ sont calculables, par exemple
par des machines de Turing ou par
récursion primitive. La relation
$\Assign{<}{N}$ est décidable.

Les modèles non standards de théories
arithmétiques comme $\Th{Q}$ et $\Th{PA}$
contiennent nécessairement des éléments
non standards. On les représente souvent
avec un domaine qui contient~$\Nat$
et d'autres éléments. Cependant, toute
!!{structure} !!{enumerable} infinie
peut être transportée sur n'importe
quel ensemble !!{denumerable}, notamment
sur~$\Nat$.\footnote{Précision éditoriale : le transport sur $\Nat$ exige un domaine infini dénombrable. La source dit seulement « dénombrable » ; une structure finie ne peut être transportée par une bijection sur $\Nat$.} Il existe donc aussi des
modèles non standards de domaine~$\Nat$.
Dans un tel modèle~$\Struct{M}$, certains
entiers naturels jouent un rôle différent
de leur rôle habituel : il existe un~$k$
qui n'est égal à
$\Value{\num{n}}{M}$ pour aucun
$n\in\Nat$.
\end{explain}

\begin{defn}
Une !!{structure}~$\Struct{M}$
pour~$\Lang{L_A}$ est \emph{calculable}
si $\Domain{M}=\Nat$, si les fonctions
$\Assign{\prime}{M}$,
$\Assign{+}{M}$ et
$\Assign{\times}{M}$ sont calculables,
et si la relation $\Assign{<}{M}$
est décidable.
\end{defn}

\begin{ex}
\ollabel{ex:comp-model-q}
Reprenons la !!{structure}~$\Struct{K}$
de l'\olref[mod][mar][mdq]{ex:model-K-of-Q}.
Son domaine était
$\Domain{K}=\Nat\cup\{a\}$,
avec les interprétations
\begin{align*}
  \Assign{\Obj{0}}{K} & = 0\\
  \Assign{\prime}{K}(x) & =
  \begin{cases}
    x+1 & \text{si $x\in\Nat$}\\
    a & \text{si $x=a$}
  \end{cases}\\
  \Assign{+}{K}(x,y) & =
  \begin{cases}
    x+y & \text{si $x,y\in\Nat$}\\
    a & \text{sinon}
  \end{cases}\\
  \Assign{\times}{K}(x,y) & =
  \begin{cases}
    xy & \text{si $x,y\in\Nat$}\\
    0 & \text{si $x=0$ ou $y=0$}\\
    a & \text{sinon}
  \end{cases}\\
  \Assign{<}{K} & =
  \Setabs{\tuple{x,y}}{x,y\in\Nat
    \text{ et }x<y}
  \cup
  \Setabs{\tuple{x,a}}{x\in\Domain{K}}.
\end{align*}
L'ensemble $\Domain{K}$ est
!!{denumerable} et peut donc être mis
en bijection avec~$\Nat$. Par exemple,
$g\colon\Nat\to\Domain{K}$, défini par
$g(0)=a$ et $g(n)=n-1$ pour $n>0$,
est une bijection.\footnote{Corrections éditoriales : la source propose $g(n)=n+1$ pour $n>0$, qui omet $0$ et $1$ du domaine cible ; $g(n)=n-1$ correspond aux opérations affichées. Dans la définition répétée de l'ordre sur $\Struct{K}$, nous rétablissons aussi la variable $x$ à la place de $n$ dans la condition d'appartenance.} Transportons les
interprétations de~$\Struct{K}$ par~$g$
pour obtenir une !!{structure}~$\Struct{K'}$
de domaine~$\Nat$, isomorphe
à~$\Struct{K}$ et donc modèle de~$\Th{Q}$.
Les fonctions et relations de~$\Lang{L_A}$
y ont des interprétations différentes,
puisque les éléments de~$\Nat$ n'y jouent
plus tous leur rôle standard.

En particulier, $0$ joue désormais
le rôle de~$a$, tandis que $1$ joue
le rôle du plus petit entier standard.
Ainsi $\Assign{\Obj{0}}{K'}=1$.
La fonction successeur satisfait
$\Assign{\prime}{K'}(0)=0$
et $\Assign{\prime}{K'}(n)=n+1$
pour $n>0$. Pour l'addition et
la multiplication, on obtient
\begin{align*}
\Assign{+}{K'}(x,y) & =
  \begin{cases}
    x+y-1 & \text{si $x,y>0$}\\
    0 & \text{sinon}
  \end{cases}\\
\Assign{\times}{K'}(x,y) & =
  \begin{cases}
    1 & \text{si $x=1$ ou $y=1$}\\
    xy-x-y+2 & \text{si $x,y>1$}\\
    0 & \text{sinon}.
  \end{cases}
\end{align*}
Enfin,
$\tuple{x,y}\in\Assign{<}{K'}$
si et seulement si
$(x>0\land y>0\land x<y)$
ou $y=0$. En particulier, $0<0$
est vrai \emph{dans}~$\Struct{K'}$,
car $a<a$ l'était dans~$\Struct{K}$.

Toutes ces fonctions sont calculables
sur~$\Nat$, et $\Assign{<}{K'}$
est une relation décidable. Ce ne
sont cependant pas les opérations
usuelles de~$\Nat$, ni son ordre usuel.
\end{ex}

\begin{prob}
Construire une !!{structure}~$\Struct{L'}$
de domaine~$\Nat$ qui soit isomorphe
à la !!{structure}~$\Struct{L}$ de
l'\olref[mod][mar][mdq]{ex:model-L-of-Q}.
\end{prob}

\begin{explain}
L'\olref{ex:comp-model-q} montre que
$\Th{Q}$ possède des modèles non standards
calculables de domaine~$\Nat$.
Le résultat suivant interdit cela
pour~$\Th{PA}$, donc aussi
pour~$\Th{TA}$.
\end{explain}

\begin{thm}[Théorème de Tennenbaum]
Tout modèle calculable de~$\Th{PA}$
est isomorphe au modèle standard
$\Struct{N}$.\footnote{Précision éditoriale : « le seul » dans la source est compris à isomorphisme près. Une renumérotation calculable du modèle standard donne une structure calculable distincte, mais isomorphe à $\Struct{N}$.}
\end{thm}

\end{document}
